\paragraph{Original introduction: proof-complexity consequences and scope}
A lower bound on the tree-like communication complexity of $\operatorname{Search}(\phi)$ implies a lower bound on the size of tree-like refutations in the corresponding proof system, and on the size of dag-like refutations using small space (size-space tradeoffs [9,12]). The usual translation from tree-like refutations of size $S$, or small-space realizations of dag-like refutations of size $S$, gives communication cost $O(\log S\log\log S\cdot c)$, with $c$ the communication cost of a proof line; sometimes it improves to $O(\log S\cdot c)$.
Randomized two-party protocols correspond to $\mathrm{Res}(\oplus)$ [15], whose lines are disjunctions of linear equations over $\mathbb F_2$. Real $k$-party NOF protocols correspond to semantic $\mathrm{Th}(k-1)$ [1], whose lines are integer-coefficient polynomial inequalities of degree at most $k-1$ on Boolean variables. Randomized $k$-party NOF protocols correspond to semantic $\mathrm T^{\mathrm{cc}}(k,c)$, whose lines are arbitrary predicates of randomized NOF communication cost at most $c$. For small $c$ this simulates $\mathrm{Th}(k-1)$ and $\mathrm{Res}(\mathrm{PC}_{k-1})$; the latter's lines are disjunctions of polynomial equations over $\mathbb F_2$ of bounded degree.
Theorem4 yields the lower bound $\exp(2^{\Omega(n/k)})$ on the size of tree-like $\mathrm T^{\mathrm{cc}}(k,c)$ (and $\mathrm{Th}(k-1)$) refutations of $\mathrm{BPHP}_{2^n}^{2^n+2^{n+k-\lfloor n/k\rfloor}}$.
Hrubes and Pudlak [10] obtained lower bounds for dag-like two-party real protocols with the same variable partition and arbitrary $m>2^\ell$; the results are formally incomparable because the present randomized model is stronger and multiparty, while [10] covers dag-like protocols and arbitrary weak binary PHP.
The randomized communication lower bound is nontrivial for $k\le\log^{1-\varepsilon}n$ with $\varepsilon>0$, yielding superpolynomial tree-like $\mathrm{Th}(k)$ lower bounds for $\mathrm{BPHP}_n^m$. Proposition6 gives short tree-like $\mathrm{Th}(\log n)$ refutations. Together with [10] it implies tree-like $\mathrm{Th}(\log n)$ cannot be simulated by dag-like $\mathrm{Th}(1)$.
Theorem4 and Proposition6 imply that $\mathrm{BPHP}_{2^\ell}^{2^\ell+1}$ separates tree-like $\mathrm{Th}(\log n)$ from tree-like $\mathrm{Th}(k)$ for $k\le\log^{1-\varepsilon}n$. Here $n=(2^\ell+1)\ell$ is the number of variables; Proposition6 gives polynomial size, while Theorem4 gives at least $\exp(\exp(\Omega(\log^\varepsilon n)))$, which is superpolynomial.
The abstract states rank lower bounds among the general consequences of randomized NOF search bounds; these are contextual consequences rather than a separately counted claim.
